% ============================================================
% IN2 - MODELLI DI CALCOLO, A.A. 2006-2007, M. PEDICINI
% ESONERO DEL 8/11/2006
% ============================================================

% ------------------------------------------------------------
% ESERCIZIO 1
% ------------------------------------------------------------
\begin{ex}
Introduciamo un nuovo modello di calcolo: l'automa a pila. Un PDFA (push-down finite automata) è una variante dell'automa a stati finiti, ed è definito nel modo seguente.

DEFINIZIONE: Un PDFA di alfabeto $A$, alfabeto della pila $B$ e insieme degli stati $Q$ è definito come una quadrupla $(T,q_0,Z_0,F)$ dove
\[
T:Q\times(A\cup\{\varepsilon\})\times B\to Q\times B^*,
\]
dove $\varepsilon$ rappresenta la parola vuota, $F:Q\to\{0,1\}$, $q_0\in Q$ è uno stato particolare detto stato iniziale, $Z_0\in B$ è il carattere con cui è inizializzata la pila.

DESCRIZIONE DELLA COMPUTAZIONE ASSOCIATA: La pila contiene ad ogni passo di computazione una parola $w\in B^*$; la funzione di transizione agisce analogamente al caso degli automi finiti in base al carattere letto nella parola di input (eventualmente $\varepsilon$), allo stato corrente e inoltre in funzione dell'ultimo carattere della pila. La transizione determinata dal valore $T(q,a,b)=(q',w')$ consiste nell'aggiornamento dello stato corrente che da $q$ diventa $q'$ (come nel caso dell'automa finito), il puntatore alla posizione corrente viene avanzato se $a\neq\varepsilon$ mentre resta fermo se $a=\varepsilon$ e l'ultimo carattere della pila viene sostituito con la parola $w'$ (quindi se la pila è $wb$ si ha $ww'$).

È necessario richiedere che in corrispondenza del carattere letto $\varepsilon$ non sia possibile applicare nessuna transizione con carattere diverso da $\varepsilon$ (altrimenti si perderebbe il determinismo).

L'automa a pila termina la computazione quando ha letto l'ultimo carattere della parola di ingresso, ha effettuato la transizione associata e non può effettuare transizioni con carattere letto $\varepsilon$. La configurazione finale è data da $q$ e dalla pila $w$; analogamente al caso degli automi finiti, la parola di input è accettata se $F(q)=1$ e viene rifiutata se $F(q)=0$.
\renewcommand{\labelenumi}{\alph{enumi})}
\ 
\begin{enumerate}
\item Definire l'algoritmo formale associato ad un automa a pila (e quindi il concetto di decidibilità per PDFA).
\item Valutare il numero di configurazioni possibili per un PDFA in funzione della lunghezza della parola di input.
\item Dimostrare che ogni $X$ decidibile per automa finito è anche PDFA decidibile.
\item Dimostrare che l'insieme $X=\{0^p1^p\mid p\in\mathbb N\}$ è PDFA-decidibile.
\item Dimostrare che se $X$ è decidibile per PDFA allora è decidibile per macchina di Turing.
\item Fornire una maggiorazione per la complessità (rispetto alla macchina di Turing) degli insiemi PDFA-decidibili.
\end{enumerate}
\end{ex}

\begin{sol}
Convenzioni. Se $(q,\varepsilon,b)\in\mathrm{dom}\,T$ allora $(q,a,b)\notin\mathrm{dom}\,T$ per ogni $a\in A$ (determinismo), quindi le $\varepsilon$-mosse hanno la precedenza. Se la pila è vuota non c'è alcun ultimo carattere e nessuna transizione è applicabile: un automa che decide $X$ deve su ogni $w$ leggere tutto l'input, quindi non si blocca prima della fine. Per $w\in A^*$ indichiamo con $w(p)$ la $p$-esima lettera.
\renewcommand{\labelenumi}{\alph{enumi})}
\ 
\begin{enumerate}
\item Lo spazio delle configurazioni è
\[
\mathrm{Conf}_{\mathrm{PDFA}}^{A,Q}=A^*\times\mathbb N\times Q\times B^*,
\]
dove $(w,p,q,\zeta)$ indica la parola di input $w$, la posizione corrente $p\in\{1,\dots,|w|+1\}$, lo stato $q$ e la pila $\zeta$. La funzione di ingresso è
\[
\iota:A^*\to\mathrm{Conf}_{\mathrm{PDFA}}^{A,Q},\qquad w\longmapsto(w,1,q_0,Z_0).
\]
La funzione di transizione agisce sull'ultimo carattere della pila: per $\zeta\in B^*$ e $b\in B$,
\[
\tau(w,p,q,\zeta b)=
\begin{cases}
(w,p,q',\zeta w') & \text{se }(q,\varepsilon,b)\in\mathrm{dom}\,T\text{ e }T(q,\varepsilon,b)=(q',w'),\\[2pt]
(w,p+1,q',\zeta w') & \text{se }(q,\varepsilon,b)\notin\mathrm{dom}\,T,\ p\leq|w|\text{ e }T(q,w(p),b)=(q',w').
\end{cases}
\]
La funzione di uscita è definita sulle configurazioni finali, cioè quelle con $p>|w|$ in cui non si può applicare nessuna $\varepsilon$-mossa (la pila è vuota oppure $(q,\varepsilon,b)\notin\mathrm{dom}\,T$ con $b$ ultimo carattere della pila):
\[
\delta:\mathrm{Conf}_{\mathrm{PDFA}}^{A,Q}\to\{0,1\},\qquad \delta(w,p,q,\zeta)=F(q).
\]
L'algoritmo associato è la sestupla
\[
\bigl(A,\{0,1\},\mathrm{Conf}_{\mathrm{PDFA}}^{A,Q},\iota,\tau,\delta\bigr).
\]
Un insieme $X\subseteq A^*$ è PDFA-decidibile se esiste un PDFA tale che per ogni $w\in A^*$ la successione $\iota(w),\tau(\iota(w)),\tau^2(\iota(w)),\dots$ raggiunge una configurazione finale $\gamma$ con $\delta(\gamma)=1$ se $w\in X$ e $\delta(\gamma)=0$ se $w\notin X$.

\item Lo spazio $\mathrm{Conf}_{\mathrm{PDFA}}^{A,Q}$ è infinito, perché la pila non ha un limite di lunghezza, e un PDFA qualsiasi può non terminare (ad esempio con $T(q,\varepsilon,b)=(q,bb)$ si continua a riempire la pila senza leggere nulla). Consideriamo quindi un PDFA $M$ che decide $X$, cioè che termina su ogni input, e contiamo le configurazioni che visita. Sia $c$ la massima lunghezza di una parola $w'$ che compare come immagine di $T$.

Vogliamo mostrare che su una parola $w$ con $|w|=n$ le configurazioni visitate sono al più $K(n+1)$, con $K$ costante che dipende solo da $M$. Poiché $M$ è deterministico e termina, le configurazioni visitate sono a due a due distinte, quindi basta contare i passi.

Passi di lettura. Ce ne sono esattamente $n$, e ognuno inserisce nella pila al più $c$ caratteri.

Passi di $\varepsilon$-mossa. Sono quelli difficili da contare, perché non leggono input. L'osservazione chiave è che una $\varepsilon$-mossa guarda solo l'ultimo carattere della pila. Se la pila è $\zeta b$ e lo stato è $q$, le $\varepsilon$-mosse successive, finché non si tocca la parte $\zeta$, dipendono solo dalla coppia $(q,b)$. Per ogni coppia $(q,b)$ ci sono quindi tre esiti possibili:
\ 
\begin{enumerate}
\item il carattere $b$ viene eliminato: dopo al più $S$ passi la pila è $\zeta$;
\item le $\varepsilon$-mosse si fermano senza toccare $\zeta$: dopo al più $S$ passi la pila è $\zeta V$ con $|V|\leq L$ e non ci sono più $\varepsilon$-mosse applicabili;
\item le $\varepsilon$-mosse non finiscono mai.
\end{enumerate}
L'esito (3) non può capitare, perché $M$ termina. Le coppie $(q,b)$ sono in numero finito, quindi $S$ e $L$ sono costanti.

Conteggio. Dopo ogni lettura, e all'inizio, c'è una fase di $\varepsilon$-mosse: sono $n+1$ fasi. Ogni fase è una successione di esiti di tipo (1) seguita al più da un esito di tipo (2).
\ 
\begin{enumerate}
\item Gli esiti di tipo (2) sono al più $n+1$, uno per fase, e ciascuno costa al più $S$ passi.
\item Ogni esito di tipo (1) elimina un carattere che era stato inserito prima. I caratteri inseriti in totale sono al più $1+nc+(n+1)L$ (quello iniziale, quelli delle letture, quelli degli esiti di tipo (2)), quindi gli esiti di tipo (1) sono al più tanti, e ciascuno costa al più $S$ passi.
\end{enumerate}
In totale
\[
\#\text{passi}\leq n+S\bigl(1+nc+(n+1)L+(n+1)\bigr)\leq K(n+1).
\]
Dunque il numero di configurazioni visitate è $O(n)$, e anche la lunghezza della pila è $O(n)$ (al più $1+cK(n+1)$).

\item Sia $X$ decidibile con un DFA $(Q,A,\delta,q_0,F)$. Poniamo $B=\{Z_0\}$ e definiamo $T(q,a,Z_0)=(\delta(q,a),Z_0)$ per ogni $q\in Q$ e $a\in A$, lasciando $T$ non definita sulle terne con $\varepsilon$; la $F$ è la stessa. La pila contiene sempre $Z_0$, non ci sono $\varepsilon$-mosse e ad ogni passo si legge una lettera: il PDFA si comporta esattamente come il DFA, quindi decide $X$.

\item Usiamo $B=\{Z_0,0\}$. Le etichette $a,b/w'$ nel disegno indicano $T(q,a,b)=(q',w')$, mentre $a$ e $b$ nel cappio di $q_r$ variano rispettivamente in $\{0,1\}$ e in $\{Z_0,0\}$. Gli stati finali sono $F(s)=F(q_a)=1$ e $F(q_0)=F(q_1)=F(q_r)=0$.
\[
\begin{tikzpicture}[
    scale=0.8, transform shape,
    >=stealth', shorten >=1pt, auto,
    every state/.style={circle, draw, fill=white, minimum size=0.8cm, inner sep=1pt}
]
% Stati
\node[state, accepting] (s)  at (0,0)    {$s$};
\node[state]            (q0) at (4,0)    {$q_0$};
\node[state]            (q1) at (8,0)    {$q_1$};
\node[state, accepting] (qa) at (12,0)   {$q_a$};
\node[state]            (qr) at (6,-3.5) {$q_r$};

% Freccia d'ingresso
\draw[->] ([shift={(-0.7,0.7)}]s.center) -- (s);

% Transizioni lungo la riga principale
\draw[->] (s)  -- node {$0,Z_0/Z_0\,0$} (q0);
\draw[->] (q0) -- node {$1,0/\varepsilon$} (q1);
\draw[->] (q1) -- node {$\varepsilon,Z_0/Z_0$} (qa);

% Cappi
\draw[->] (q0) to[loop above, looseness=8] node {$0,0/0\,0$} (q0);
\draw[->] (q1) to[loop above, looseness=8] node {$1,0/\varepsilon$} (q1);
\draw[->] (qr) to[loop below, looseness=8] node {$a,b/b$} (qr);

% Transizioni verso lo stato di rifiuto
\draw[->] (s)  -- node[swap, pos=0.4] {$1,Z_0/Z_0$} (qr);
\draw[->] (q1) -- node[swap, pos=0.4] {$0,0/0$} (qr);
\draw[->] (qa) -- node[pos=0.4] {$a,Z_0/Z_0$} (qr);
\end{tikzpicture}
\]
Il carattere $Z_0$ non viene mai eliminato, quindi la pila non è mai vuota. Il test di fine input non è visibile a $T$, perciò si usa $F$ insieme alla $\varepsilon$-mossa di $q_1$, che scatta appena la pila contiene solo $Z_0$. Sia $w\in\{0,1\}^*$.
\ 
\begin{enumerate}
\item Se $w=\varepsilon$ il calcolo termina in $s$ e $F(s)=1$: $w$ è accettata.
\item Se $w=0^p1^q$ con $p\geq1$, dopo $0^p$ si è in $q_0$ con pila $Z_0\,0^p$. Le $q$ cifre $1$ eliminano $q$ zeri: se $q<p$ si termina in $q_1$ con pila $Z_0\,0^{p-q}$ e $F(q_1)=0$, se $q=p$ la $\varepsilon$-mossa porta in $q_a$ e $F(q_a)=1$, se $q>p$ la cifra $1$ successiva porta da $q_a$ in $q_r$. Per $q=0$ si termina in $q_0$ con $F(q_0)=0$.
\item In tutti gli altri casi (la parola inizia per $1$, oppure una cifra $0$ segue una cifra $1$) si finisce in $q_r$, che non è finale.
\end{enumerate}
Quindi $w$ è accettata se e solo se $w=0^p1^p$ per qualche $p\geq0$.

\item Sia $M=(T,q_0,Z_0,F)$ un PDFA che decide $X$. Costruiamo una macchina di Turing $\mathcal M$ a due nastri, con alfabeto $A\cup B\cup\{\square\}$. Sul primo nastro c'è la parola di input, con la testina sulla posizione corrente $p$; sul secondo nastro c'è la pila, con la testina sull'ultimo carattere, e all'inizio vi è scritto $Z_0$. Lo stato di $\mathcal M$ contiene lo stato corrente $q$ di $M$. Ad ogni passo, siano $a$ il carattere sotto la testina del primo nastro e $b$ quello sotto la testina del secondo ($\square$ se la pila è vuota).
\ 
\begin{enumerate}
\item Se $(q,\varepsilon,b)\in\mathrm{dom}\,T$ con $T(q,\varepsilon,b)=(q',w')$, si sostituisce $b$ con $w'$ sul secondo nastro, lasciando ferma la testina del primo.
\item Altrimenti, se $a\in A$ e $T(q,a,b)=(q',w')$, si sostituisce $b$ con $w'$ sul secondo nastro e si sposta a destra la testina del primo.
\item Altrimenti ($a=\square$ e nessuna $\varepsilon$-mossa), $\mathcal M$ si ferma accettando se $F(q)=1$ e rifiutando se $F(q)=0$.
\end{enumerate}
Sostituire $b$ con $w'=u_1\cdots u_m$ vuol dire scrivere $u_1$ al posto di $b$ e poi $u_2,\dots,u_m$ spostandosi a destra, oppure, se $w'=\varepsilon$, scrivere $\square$ e spostarsi a sinistra. Ogni passo di $M$ richiede quindi al più $c$ passi di $\mathcal M$, e $\mathcal M$ termina perché $M$ termina. Quindi $\mathcal M$ decide $X$.

\item Per il lemma del punto (b) $M$ fa al più $K(n+1)$ passi, quindi la macchina $\mathcal M$ del punto (e) lavora in tempo al più $cK(n+1)$ e
\[
X\in\mathrm{DTIME}_2(c'n).
\]
Per lo spazio, sul primo nastro si usano $n+1$ celle, mentre sul secondo se ne usano al più $1+cK(n+1)$: lo spazio è $O(n)$ su entrambi i nastri, e non $O(1)$ sul secondo (per $0^p1^p$ la pila contiene $p$ zeri). Passando a una macchina a un solo nastro con la simulazione standard si ottiene $X\in\mathrm{DTIME}_1(c''n^2)$.
\end{enumerate}
\end{sol}
\newpage

% ------------------------------------------------------------
% ESERCIZIO 2
% ------------------------------------------------------------
\begin{ex}
(facoltativo) Con le notazioni dell'esercizio 1, si consideri la seguente variante PDFA' di automa a pila, che ridefinisce il solo concetto di parola accettata dall'automa nel modo seguente: l'automa a pila termina la computazione quando ha letto l'ultimo carattere della parola di ingresso ed ha effettuato la transizione associata e non può effettuare transizioni con carattere letto $\varepsilon$, terminando con uno stato $q$ e con la pila $w$: in questo caso la parola di input è accettata se $F(q)=1$ e $w$ è la parola vuota e viene rifiutata se $F(q)=0$ oppure se $w$ non è vuota.

Provare che $X$ è PDFA decidibile se e solo se $X$ è PDFA' decidibile.
\end{ex}

\begin{sol}
Vale l'implicazione \glqq PDFA'-decidibile $\Rightarrow$ PDFA-decidibile\grqq, mentre l'implicazione opposta è falsa.
\ 
\begin{enumerate}
\item Sia $M=(T,q_0,Z_0,F)$ un PDFA' che decide $X$. Costruiamo un PDFA $M'=(T',q_0',Z_0',F')$ con stati $Q'=Q\cup\{q_0',q_a,q_r\}$ e alfabeto della pila $B'=B\cup\{Z_0'\}$, dove $q_0'$, $q_a$, $q_r$ e $Z_0'$ sono nuovi:
\ 
\begin{enumerate}
\item $T'(q_0',\varepsilon,Z_0')=(q_0,Z_0'Z_0)$;
\item $T'(q,a,b)=T(q,a,b)$ per $q\in Q$, $a\in A\cup\{\varepsilon\}$, $b\in B$;
\item per $q\in Q$, $T'(q,\varepsilon,Z_0')=(q_a,Z_0')$ se $F(q)=1$ e $T'(q,\varepsilon,Z_0')=(q_r,Z_0')$ se $F(q)=0$;
\item $T'(q_a,a,Z_0')=T'(q_r,a,Z_0')=(q_r,Z_0')$ per $a\in A$;
\item $F'(q_a)=1$ e $F'(q)=0$ per ogni altro stato.
\end{enumerate}
Il determinismo è rispettato: su $Z_0'$ sono definite solo $\varepsilon$-mosse (oppure, per $q_a$ e $q_r$, solo mosse con lettura).

Sull'input $w$, $M'$ fa prima l'$\varepsilon$-mossa che porta in $(q_0,Z_0'Z_0)$ e poi simula $M$ passo per passo, con la pila di $M$ sopra $Z_0'$. Quando $M$ termina in uno stato $q$ con pila $\zeta$ ci sono due casi.
\ 
\begin{enumerate}
\item Se $\zeta\neq\varepsilon$, $M'$ ha in cima un carattere di $B$, non ha $\varepsilon$-mosse e ha letto tutto l'input: termina in $q\in Q$ con $F'(q)=0$, e anche $M$ rifiuta.
\item Se $\zeta=\varepsilon$, $M'$ ha $Z_0'$ in cima e l'$\varepsilon$-mossa porta in $q_a$ se $F(q)=1$ e in $q_r$ se $F(q)=0$, dove $M'$ termina: $M'$ accetta se e solo se $M$ accetta.
\end{enumerate}
Se la pila di $M$ si svuotasse prima della fine dell'input, $M$ si bloccherebbe, ma questo non accade perché $M$ decide $X$. Quindi $M'$ decide $X$.

\item Sia $A=\{0\}$ e $X=\{0\}^*$. $X$ è decidibile con un DFA a un solo stato finale, quindi è PDFA-decidibile per il punto (c). Supponiamo che un PDFA' $M$ decida $X$. Poiché $\varepsilon\in X$, sull'input $\varepsilon$ il calcolo termina con pila vuota (dopo le $\varepsilon$-mosse iniziali) e $F(q)=1$. Sull'input $0$ le $\varepsilon$-mosse iniziali sono le stesse, perché non dipendono dall'input, e portano alla stessa configurazione con pila vuota, ma con il carattere $0$ ancora da leggere: nessuna transizione è applicabile, quindi $M$ non decide se $0\in X$. Dunque $X$ non è PDFA'-decidibile.
\end{enumerate}
\begin{rema}
Il motivo è che $T$ non vede la fine dell'input, quindi il PDFA' non può svuotare la pila \glqq alla fine\grqq: dopo ogni prefisso accettato la pila dovrebbe essere già vuota. Lo stesso accade se \glqq pila vuota\grqq{} si intende \glqq pila uguale a $Z_0$\grqq. Ad esempio $X=\{0^p1^q\mid q\leq p\}$ è PDFA-decidibile (si conta con la pila), ma non PDFA'-decidibile. Infatti ogni $0^p$ è in $X$, quindi dopo $0^p$ la configurazione è $(q_p,Z_0)$ con $q_p$ in un insieme finito; per $p<p'$ con $q_p=q_{p'}$ le parole $0^p1^{p'}\notin X$ e $0^{p'}1^{p'}\in X$ avrebbero lo stesso esito, assurdo.
\end{rema}
\end{sol}
\newpage

% ------------------------------------------------------------
% ESERCIZIO 3
% ------------------------------------------------------------
\begin{ex}
Si consideri la seguente funzione $T:\mathbb N\to\mathbb N$, $b\in\mathbb N$:
\[
T(n)=
\begin{cases}
b & \text{se }n<3,\\
3T(n/3)+bn & \text{se }n\geq3.
\end{cases}
\]
\renewcommand{\labelenumi}{\alph{enumi})}
\ 
\begin{enumerate}
\item Mostrare che $T(n)$ è una funzione di complessità;
\item Mostrare che $\mathrm{DTIME}(T(n))=\mathrm{DTIME}(n\log(n))$.
\end{enumerate}
\end{ex}

\begin{sol}
Qui $n/3$ indica $\lceil n/3\rceil$ (con la parte intera inferiore i calcoli sono identici), e una funzione di complessità è una funzione $T:\mathbb N\to\mathbb N$ non decrescente con $T(n)\geq n$.
\renewcommand{\labelenumi}{\alph{enumi})}
\ 
\begin{enumerate}
\item Mostriamo prima che $T$ è non decrescente, per induzione forte su $n$, cioè che $T(n)\leq T(n+1)$ per ogni $n$.
\ 
\begin{enumerate}
\item Per $n+1\leq2$ si ha $T(n)=T(n+1)=b$, e per $n=2$ si ha $T(3)=3T(1)+3b=6b\geq b=T(2)$.
\item Sia $n\geq3$ e supponiamo $T$ non decrescente fino a $n$. Poiché $\lceil n/3\rceil\leq\lceil(n+1)/3\rceil\leq n$, l'ipotesi induttiva dà
\[
T(n+1)=3T\bigl(\lceil\tfrac{n+1}3\rceil\bigr)+b(n+1)\geq3T\bigl(\lceil\tfrac n3\rceil\bigr)+bn=T(n).
\]
\end{enumerate}
Mostriamo ora che $T(n)\geq n$. Se $b=0$ allora $T\equiv0$ e $T$ non è di complessità. Sia $b\geq1$. Per $n\geq3$ si ha $T(n)\geq bn\geq n$, e $T(1)=b\geq1$. Per $n=2$ si ha $T(2)=b\geq2$ se $b\geq2$; per $b=1$ la disuguaglianza fallisce solo per $n=2$ ($T(2)=1<2$), cosa irrilevante per l'andamento asintotico. Dunque $T$ è di complessità per $b\geq2$ (e, a meno del valore in $n=2$, per $b=1$).

\item Sia $b\geq1$. Per $n=3^k$, per induzione su $k$ si ha $T(3^k)=b\,3^k(k+1)$: per $k=0$ è $T(1)=b$ e
\[
T(3^{k+1})=3T(3^k)+b\,3^{k+1}=b\,3^{k+1}(k+1)+b\,3^{k+1}=b\,3^{k+1}(k+2).
\]
In particolare, lungo $n=3^k$,
\[
\frac{T(n)}{n\log n}=\frac{b\,3^k(k+1)}{3^k\,k\log3}\longrightarrow\frac b{\log3}.
\]
Per ottenere il confronto per ogni $n$, sia $3^k\leq n<3^{k+1}$ con $k\geq1$. Poiché $T$ è non decrescente,
\[
b\,3^k(k+1)\leq T(n)\leq b\,3^{k+1}(k+2),
\qquad
3^k\,k\log3\leq n\log n\leq3^{k+1}(k+1)\log3,
\]
da cui
\[
\frac{b}{3\log3}\leq\frac{T(n)}{n\log n}\leq\frac{3b(k+2)}{k\log3}\leq\frac{9b}{\log3}.
\]
Esistono dunque $c_1,c_2>0$ con $c_1\,n\log n\leq T(n)\leq c_2\,n\log n$ per ogni $n\geq3$. Poiché $\mathrm{DTIME}(cf)=\mathrm{DTIME}(f)$ per ogni costante $c$ (accelerazione lineare) e $\mathrm{DTIME}$ è monotona rispetto alla funzione, da $T\leq c_2\,n\log n$ segue $\mathrm{DTIME}(T(n))\subseteq\mathrm{DTIME}(n\log n)$ e da $n\log n\leq T(n)/c_1$ segue l'inclusione opposta. Quindi
\[
\mathrm{DTIME}(T(n))=\mathrm{DTIME}(n\log n).
\]
\end{enumerate}
\end{sol}
\newpage